\section{Protocolos experimentales}
\label{sec:protocolos}

Esta sección fija los experimentos antes de presentar sus resultados. La
separación es importante porque compartir una familia de ecuaciones o un mismo
conjunto de parámetros no convierte dos trayectorias en la misma
reproducción. En particular, la reconstrucción de Danca y la ruta no suave
propuesta usan los parámetros de \cref{eq:parametros_no_suave}, pero emplean
semillas y procedimientos de continuación diferentes. De manera análoga, la
reconstrucción de Wu y el caso \texttt{c590} pertenecen a
\cref{eq:chua_arctan}, aunque usan parámetros, condiciones iniciales e
integradores distintos.

\begin{table}[H]
\centering
\caption{Carriles experimentales y función de cada uno en el reporte.}
\label{tab:carriles_experimentales}
\begin{tabularx}{\linewidth}{L{0.23\linewidth}L{0.25\linewidth}Y}
\toprule
Carril & Contrato dinámico & Propósito\\
\midrule
Leonov--Kuznetsov, \(q=1\)
& Sistema entero y control EFORK--3 en \(q=1\)
& Recuperar la construcción armónica y el procedimiento entero\\
Danca, \(q=0.9998\)
& Caputo ABM con memoria completa
& Reconstruir los datos publicados y evaluar una semilla diagnóstica
declarada\\
Wu, \(q=0.99\)
& Recurrencia ADM local de orden cuatro
& Reintegrar las dos condiciones iniciales publicadas\\
No suave propuesto
& Función descriptiva sesgada, continuación Caputo y ABM
& Localizar y sondear el conjunto atrayente obtenido con la nueva semilla\\
\texttt{c590}
& Búsqueda RK4 en \(q=1\), refinamiento Caputo ABM
& Seleccionar parámetros, refinar la semilla y sondear el conjunto final\\
\bottomrule
\end{tabularx}
\end{table}

\subsection{Integración causal mediante ABM--PECE}
\label{subsec:protocolo_abm}

Los problemas de Caputo de las rutas propuestas y de la reconstrucción de
Danca se integran con el predictor--corrector
Adams--Bashforth--Moulton. Para
\[
{}^{\mathrm C}D_{t_0}^{q}X(t)=F(t,X(t)),\qquad X(t_0)=X_0,
\]
se discretiza la ecuación de Volterra de \cref{eq:volterra} en
\(t_n=t_0+nh\). Si \(F_j=F(t_j,X_j)\), el predictor es
\begin{equation}
X_{n+1}^{\mathrm P}
=X_0+\frac{h^q}{\Gamma(q+1)}
\sum_{j=0}^{n}b_{j,n+1}F_j,
\qquad
b_{j,n+1}=(n+1-j)^q-(n-j)^q .
\label{eq:abm_predictor_protocolo}
\end{equation}
El campo se evalúa una vez en el estado predicho y se aplica el corrector
\begin{equation}
X_{n+1}
=X_0+\frac{h^q}{\Gamma(q+2)}
\left[
F(t_{n+1},X_{n+1}^{\mathrm P})
+\sum_{j=0}^{n}a_{j,n+1}F_j
\right],
\label{eq:abm_corrector_protocolo}
\end{equation}
con
\begin{equation}
a_{j,n+1}=
\begin{cases}
n^{q+1}-(n-q)(n+1)^q, & j=0,\\
(n-j+2)^{q+1}-2(n-j+1)^{q+1}+(n-j)^{q+1},
&1\leq j\leq n.
\end{cases}
\label{eq:abm_pesos_protocolo}
\end{equation}

En cada paso se realizan, en este orden, las siguientes operaciones:

\begin{enumerate}
  \item se conservan todos los pares \((X_j,F_j)\) calculados desde \(t_0\);
  \item se forma la suma histórica del predictor
  \cref{eq:abm_predictor_protocolo};
  \item se evalúa el campo en \(X_{n+1}^{\mathrm P}\);
  \item se forma la suma histórica del corrector
  \cref{eq:abm_corrector_protocolo};
  \item se agregan \(X_{n+1}\) y \(F_{n+1}\) a la historia.
\end{enumerate}

Ninguno de los términos \(j=0,\ldots,n\) se elimina. Por tanto, el costo
directo es \(O(N^2)\) y el almacenamiento es \(O(N)\) para \(N\) pasos
\cite{DiethelmFordFreed2002}. Cada sonda de cuenca constituye un problema de
valor inicial independiente con el mismo instante \(t_0\) y su propia memoria.
En la continuación no suave, en cambio, los 21 niveles comparten una sola
historia; sólo al terminar la continuación se inicia el problema autónomo que
define la nube objetivo.

\subsection{Métrica de contacto y clasificación de destinos}
\label{subsec:clasificador}

Sea \(\mathcal C_h=\{A_k\}\) la nube postransitoria del conjunto objetivo y sea
\(Y=\{Y_j\}\) el segmento postransitorio de una sonda. Después del submuestreo
declarado, la distancia usada para decidir un contacto es
\begin{equation}
D_{90}(Y,\mathcal C_h)
=Q_{0.90}\left(
\left\{\min_k\|Y_j-A_k\|_2\right\}_j
\right).
\label{eq:d90_protocolo}
\end{equation}
Se registra un contacto cuando
\begin{equation}
D_{90}(Y,\mathcal C_h)\leq\delta .
\label{eq:contacto_protocolo}
\end{equation}
El percentil 90 evita que unos pocos puntos cercanos determinen por sí solos
la clasificación. Antes de ejecutar los sondeos se fijan la nube, el
submuestreo y \(\delta\); el umbral no se ajusta después de observar los
resultados.

Cada trayectoria se asigna a una de cuatro clases:

\begin{enumerate}
  \item \emph{objetivo}, si satisface \cref{eq:contacto_protocolo};
  \item \emph{equilibrio}, si satisface el criterio estacionario declarado;
  \item \emph{divergencia}, si alcanza el umbral de crecimiento;
  \item \emph{otro destino}, para las trayectorias finitas restantes.
\end{enumerate}

\subsection{Protocolos de reconstrucción bibliográfica}
\label{subsec:protocolos_bibliograficos}

\subsubsection{Control entero de Leonov--Kuznetsov: carril histórico EFORK--3}

Se usa \cref{eq:chua_entero,eq:parametros_enteros}. El cierre armónico
selecciona la primera rama positiva y construye la condición modal. El control
dinámico almacenado utiliza EFORK--3 evaluado en \(q=1\), con paso \(h=0.01\).
La prueba de vecindades emplea
\[
r\in\{10^{-5},3{\times}10^{-5},10^{-4},3{\times}10^{-4},
10^{-3},3{\times}10^{-3},10^{-2}\}
\]
y 24 condiciones por radio alrededor de cada uno de los tres equilibrios. El
conteo total es
\[
3\ \text{equilibrios}\times7\ \text{radios}\times24\ \text{muestras}
=504.
\]
Este expediente usa \(T=320\), \(T_{\mathrm{burn}}=80\) y el clasificador
histórico de huellas de sección contra 89 puntos de referencia. Sus
coordenadas de sondeo no son el banco compartido de la comparación
Python--Julia.
Este carril comprueba la recuperación numérica del caso \(q=1\); no suministra
la semilla de ninguno de los dos casos fraccionarios propuestos.

\subsubsection{Reconstrucción de Danca}

Se conservan la ecuación \cref{eq:chua_no_suave}, los parámetros
\cref{eq:parametros_no_suave} y el orden \(q=0.9998\). La condición inicial
diagnóstica se obtiene al aplicar el cierre armónico entero a esos parámetros:
\begin{equation}
X_0^{\mathrm D}
=(5.8561450863,\ 0.3693315782,\ -8.3665361683)^{\T}.
\label{eq:semilla_danca_diagnostica}
\end{equation}
Esta condición se declara como una reconstrucción y no como una condición
inicial publicada por Danca. La trayectoria se calcula mediante Caputo
ABM--PECE con memoria completa,
\[
h=0.01,\qquad T=500,\qquad T_{\mathrm{burn}}=250.
\]
La clasificación postransitoria utiliza espectro de potencia y prueba 0--1.
El protocolo de vecindad descrito en la referencia se conserva como control:
200 condiciones con \(\delta=0.01\) alrededor de un equilibrio externo y la
simetría para el equilibrio opuesto \cite{Danca2017}.

\subsubsection{Reconstrucción de Wu}

Se emplean \cref{eq:chua_arctan,eq:parametros_wu} y las condiciones iniciales
publicadas
\begin{equation}
X_{0,+}^{\mathrm W}
=(13.8,\ 0.7093,\ -19.8768)^{\T},
\qquad X_{0,-}^{\mathrm W}=-X_{0,+}^{\mathrm W}.
\label{eq:semillas_wu_protocolo}
\end{equation}
La recurrencia se ejecuta con el ADM local de orden cuatro usado en la
reproducción bibliográfica:
\[
q=0.99,\qquad h=0.01,\qquad N=10\,000,\qquad
T=100,\qquad T_{\mathrm{burn}}=50.
\]
Este actualizador local no acumula una historia ABM de Caputo. Por ello, sus
resultados se presentan únicamente dentro del carril de Wu y no se combinan
con \texttt{c590}.

\subsection{Ruta propuesta para el sistema no suave}
\label{subsec:protocolo_no_suave}

\subsubsection{Exploración y selección de la semilla}

La exploración considera nueve pares de pendientes:
\begin{equation}
m_1\in\{-1.1468,-1.20,-1.25\},\qquad
m_0\in\{-0.1768,-0.20,-0.24\}.
\label{eq:malla_pendientes_no_suave}
\end{equation}
Para cada par se inicializa el solucionador con cinco valores de
\(A\in[0.5,10]\), nueve de \(c\in[-8,8]\) y seis de
\(\omega\in[0.5,6]\). Por tanto,
\begin{equation}
5\times9\times6=270\ \text{inicios por par},\qquad
9\times270=2\,430\ \text{inicios totales}.
\label{eq:conteo_busqueda_no_suave}
\end{equation}

La ecuación de cierre sesgada se resuelve mediante mínimos cuadrados de región
de confianza reflectiva con
\[
A\in[10^{-6},25],\qquad c\in[-12,12],\qquad
\omega\in[0.1,8],
\]
\[
f_{\mathrm{tol}}=x_{\mathrm{tol}}=10^{-10},\qquad
g_{\mathrm{tol}}=10^{-8},\qquad N_{\mathrm{eval}}\leq300.
\]
Las integrales de la función descriptiva usan una regla trapezoidal periódica
con 8\,192 nodos. Una raíz pasa a la etapa siguiente si el solucionador
converge, \(A\geq0.5\), \(0.5\leq\omega\leq6\) y la norma del residuo es menor
que \(10^{-4}\). Las raíces se deduplican con tolerancia \(10^{-3}\) en
\((A,c,\omega)\), se conservan como máximo dos por par y, dentro del par
nominal de \cref{eq:parametros_no_suave}, se selecciona la raíz con
\(c>0.05\) y menor residuo.

\subsubsection{Continuación con una sola memoria}

La semilla seleccionada se transporta mediante
\[
\varepsilon_k=0.05k,\qquad k=0,\ldots,20.
\]
Cada uno de los 21 niveles usa \(h=0.01\), 30 unidades de transitorio y 30
unidades retenidas. El integrador no se reinicia al cambiar
\(\varepsilon_k\): las sumas de
\cref{eq:abm_predictor_protocolo,eq:abm_corrector_protocolo} mantienen todos
los estados calculados desde el inicio del nivel \(\varepsilon=0\). El estado
final en \(\varepsilon=1\) inicia después un nuevo problema autónomo con
\[
q=0.9998,\qquad h=0.01,\qquad T=300,\qquad
T_{\mathrm{burn}}=100.
\]
La atracción local del objetivo se comprueba con 12 perturbaciones en cada uno
de los radios \(10^{-6}\), \(10^{-4}\) y \(10^{-2}\).

\subsubsection{Sondeo volumétrico}

Alrededor del equilibrio \(e_i\), una condición se genera mediante
\begin{equation}
X_{0,i\ell m}
=e_i+r_\ell U_{i\ell m}^{1/3}d_{i\ell m},
\qquad
U_{i\ell m}\sim\mathcal U(0,1),\quad
d_{i\ell m}\sim\mathcal U(\mathbb S^2).
\label{eq:muestreo_bola}
\end{equation}
El factor \(U^{1/3}\) produce una distribución uniforme respecto del volumen
de la bola. El plan nominal por equilibrio es
\begin{equation}
\begin{aligned}
r_\ell\in\{&
10^{-5},3{\times}10^{-5},10^{-4},3{\times}10^{-4},
10^{-3},3{\times}10^{-3},10^{-2},\\
&0.03,0.1,0.3,1,2\},\\
N_\ell\in\{&
80,120,160,240,400,600,800,1000,1200,1200,1500,1500\}.
\end{aligned}
\label{eq:plan_bolas_no_suave}
\end{equation}
Se usa PCG64 con semilla
\[
s_{i\ell}=42+100i+10\ell .
\]
Cada sonda utiliza ABM--PECE de memoria completa con
\[
q=0.9998,\quad h=0.01,\quad T=200,\quad
T_{\mathrm{burn}}=50,\quad\delta=0.5.
\]
La convergencia temprana requiere \(t\geq25\), distancia menor que
\(10^{-3}\) a un equilibrio, norma del campo menor que \(10^{-4}\) y 50 pasos
consecutivos. La divergencia temprana requiere cinco normas consecutivas
mayores que 80 o cinco razones de crecimiento mayores que 1.25; el corte duro
es 120. Una trayectoria que completa el horizonte también se clasifica como
equilibrio si su distancia final a alguno de ellos no excede 0.5.

\subsection{Ruta propuesta para el caso \texttt{c590}}
\label{subsec:protocolo_c590}

\subsubsection{Exploración global y local en
\texorpdfstring{\(q=1\)}{q=1}}

La exploración global genera 2\,400 casos con PCG64 y semilla
\(2026062304\):
\begin{equation}
\begin{aligned}
\alpha&\sim\mathcal U(5,22),&
\beta&\sim\mathcal U(6,35),&
\gamma&=10^{\mathcal U(-3,-0.5)},\\
a_1&\sim\mathcal U(-0.1,1),&
a_2&\sim\mathcal U(-3.5,-0.4),&
\rho&\sim\mathcal U(0.4,2.8).
\end{aligned}
\label{eq:distribucion_global_c590}
\end{equation}
Cada parámetro se integra simultáneamente desde una semilla candidata y desde
el control \((10^{-3},0,0)^{\T}\), con RK4 fijo,
\[
q=1,\qquad h=0.01,\qquad T=200,\qquad T_{\mathrm{burn}}=100.
\]
La respuesta candidata debe permanecer acotada y separada de la respuesta del
control según la distancia \(D_{90}\). Entre las respuestas no triviales se
selecciona el caso de mayor puntuación exploratoria.

Alrededor de ese caso se generan 1\,000 perturbaciones con PCG64 y semilla
\(2026062305\):
\begin{equation}
\begin{aligned}
\alpha&=\alpha_b\exp\mathcal N(0,0.18),&
\beta&=\beta_b\exp\mathcal N(0,0.18),\\
\gamma&=\gamma_b\exp\mathcal N(0,0.35),&
a_1&=a_{1,b}+\mathcal N(0,0.12),\\
a_2&=a_{2,b}\exp\mathcal N(0,0.15),&
\rho&=\rho_b\exp\mathcal N(0,0.15).
\end{aligned}
\label{eq:perturbacion_local_c590}
\end{equation}
La integración local usa \(h=0.01\), \(T=240\) y
\(T_{\mathrm{burn}}=120\). Los 30 casos no periódicos con mayor puntuación se
auditan mediante RK4 variacional y ortonormalización QR, con
\[
h=0.005,\qquad T_{\mathrm{burn}}=200,\qquad
T_{\mathrm{med}}=300,
\]
y una reortogonalización cada diez pasos. Se selecciona el mayor exponente
principal entre los casos acotados, separados del control, con \(E_0\)
inestable y \(E_\pm\) estables.

\subsubsection{Refinamiento fraccionario y selección de la condición inicial}

Los órdenes
\[
q\in\{0.9995,0.9998,0.9999,0.99995\}
\]
se evalúan como problemas de Caputo independientes. Cada problema empieza en
\(t_0=0\) con la semilla entera seleccionada y usa ABM de memoria completa,
\[
h=0.005,\qquad T=300,\qquad T_{\mathrm{burn}}=150.
\]
De la trayectoria con \(q=0.9999\) se extraen 16 estados, en los índices
\[
30\,001,32\,000,\ldots,60\,000,
\]
que corresponden aproximadamente a \(t=150.005,\ldots,300\). Cada estado
inicia un nuevo problema de Caputo. La criba corta usa \(h=0.0025\),
\(T=180\) y \(T_{\mathrm{burn}}=90\); los supervivientes se verifican además
con \(h=0.005\) y \(h=0.01\). La auditoría larga usa
\[
h\in\{0.0025,0.005,0.01\},\qquad
T=300,\qquad T_{\mathrm{burn}}=150.
\]
La selección exige acotación en los tres pasos y mediana remuestreada de la
prueba 0--1 mayor que 0.8 en los dos pasos refinados
\(h=0.0025\) y \(h=0.005\).

\subsubsection{Sondeo sobre superficies esféricas}

Las condiciones iniciales son
\begin{equation}
X_{0,i\ell m}=e_i+r_\ell d_{i\ell m},
\qquad d_{i\ell m}\in\mathbb S^2,
\label{eq:muestreo_superficie_c590}
\end{equation}
donde las direcciones combinan los seis semiejes cartesianos y una
distribución determinista de Fibonacci esférica. El plan por equilibrio es
\begin{equation}
\begin{aligned}
r_\ell&\in\{10^{-5},10^{-4},10^{-3},10^{-2},0.03,0.1,0.3,1,2\},\\
N_\ell&\in\{100,200,300,400,500,600,700,800,900\}.
\end{aligned}
\label{eq:plan_superficies_c590}
\end{equation}
Cada sonda usa
\[
q=0.9999,\qquad h=0.005,\qquad
T=180,\qquad T_{\mathrm{burn}}=90.
\]
La nube objetivo se simetriza como
\(\mathcal C_h^\pm=\mathcal C_h\cup(-\mathcal C_h)\). Su calibración fija
\begin{equation}
d_{90}^{\mathrm{cal}}=0.2950508245,\qquad
\delta=2d_{90}^{\mathrm{cal}}=0.5901016489.
\label{eq:umbral_c590}
\end{equation}
La detección temprana de equilibrio se activa en \(t=45\) y requiere distancia
menor que \(10^{-6}\), norma del campo menor que \(10^{-4}\) y 50 pasos
consecutivos. El umbral de divergencia es \(\|X\|_2=120\).

\subsection{Partición local y auditoría macroscópica}
\label{subsec:detencion_primer_contacto}

Los radios se procesan en orden ascendente y se completa cada bloque alrededor
de \(E_0\), \(E_+\) y \(E_-\). La partición local comprende \(r\leq0.3\).
La auditoría macroscópica conserva además \(r=1\) y \(r=2\). Se define
\[
H_\ell=\sum_{i,m}
\mathbf 1\!\left[
D_{90}(Y_{i\ell m},\mathcal C_h)\leq\delta
\right],
\]
y se registra
\begin{equation}
\ell_\star=\min\{\ell:H_\ell>0\}.
\label{eq:primer_radio_contacto}
\end{equation}
Este índice es un estadístico del primer radio de contacto, no una regla de
truncamiento del registro canónico. Para \texttt{c590}, \(\ell_\star\)
corresponde a \(r=1\); el bloque \(r=2\) se retiene para cuantificar la
sensibilidad macroscópica y las divergencias.
